\subsection{An exact phase on the lower strip}

The lower strip contracts in the normal direction while its slow phase
moves at order $\epsilon$.  The horizontal coordinate of the
compactified chart also moves during fast contraction, so it cannot
serve as this phase.  We first project along the zero-scale fast
fibres, then correct the projection at positive scale to obtain an
exact scalar equation.  In the incoming chart put
\[
 T=(-X)^{-1},\qquad v=Y/(-X),\qquad w=v+\epsilon T=\epsilon zT,
 \qquad \frac{\dd\zeta}{\dd\xi}=-X>0.
\]
Direct substitution in the scaled original field gives
\begin{equation}\label{rw:transitions:compact-lower}
\begin{aligned}
 \dot T={}&\epsilon T[\bar B-\bar aT+(1+\epsilon\bar m)T^2]
       -(1+2\epsilon\bar m)T^2w+\bar mTw^2,\\
 \dot w={}&\epsilon^2\nu T^2+
 [-1+\epsilon\bar B+\epsilon(1+\epsilon\bar m)T^2
                  -\epsilon(\bar a+\bar c+2\nu)T]w\\
 &+[\bar c+\nu-(1+2\epsilon\bar m)T]w^2+\bar m w^3.
\end{aligned}
\end{equation}
Dots in this subsection denote $\zeta$-time. The outgoing chart is
obtained by reversing time and replacing
$(\bar a,\bar c,\nu)$ by $(-\bar a,-\bar c,-\nu)$.

Fix a small $R>0$ such that
$1\pm\bar aX+\bar BX^2>3/4$ for $0\le X\le2R$ throughout
the chosen bounded coefficient box, and put $T_R=1/R$.
All constructions below take place first on a slightly larger
$T$-interval than $[0,T_R]$.

\begin{proposition}[Lower phase]\label{rw:phase}
For each prescribed finite order $N$, there are a normal width and a
scale cutoff for which the following statements hold on common smaller
state and coefficient domains. There is an invariant graph
\begin{equation}\label{rw:transitions:weak-graph}
 w_g(T)=\epsilon^2\nu T^2H(T;\mathbf b,\epsilon),
 \qquad H(0;\mathbf b,\epsilon)=\frac1{1+\epsilon\bar B},
\end{equation}
and coordinates $(x,Z)$, with $Z=w-w_g(T)$ and $x=T$ on that graph,
in which the field is exactly
\begin{equation}\label{rw:transitions:lower-triangular}
 \dot x=\epsilon g_I(x),\qquad
 g_I(x)=xF_I(x),\qquad \dot Z=-Z\beta_I(x,Z).
\end{equation}
Their finite mixed derivatives, including the ordinary scale derivatives
used in their construction, are bounded. The normal rates satisfy
\[
 \beta_I>\beta_*,\qquad
 \beta_I+Z\partial_Z\beta_I>\beta_*,\qquad
 \beta_I(0,0)=1-\epsilon\bar B.
\]
Moreover,
\begin{equation}\label{rw:transitions:lower-factor}
\begin{aligned}
 F_I(x)={}&\bar B-\bar ax+(1+\epsilon\bar m)x^2
 -\epsilon\nu(1+2\epsilon\bar m)x^3H(x)
 +\epsilon^3\nu^2\bar m x^4H(x)^2,\\
 F_I(0)={}&\bar B,\qquad F_I'(0)=-\bar a.
\end{aligned}
\end{equation}
Thus $F_I''>1$ after the scale cutoff is reduced, and $F_I$ has a
positive margin near the fixed exit $x_R=T_R$.
In the derivative convention of this section,
$\beta_I(x,0)=1+O_{C^N}(\epsilon)$.
Ordinary differentiation with respect to $\epsilon$ may remove the
factor $\epsilon$ in this last estimate.
\end{proposition}

\begin{proof}
At $\epsilon=0$ the field is
\[
 \dot T=Tw(\bar m w-T),\qquad
 \dot w=w[-1+(\bar c+\nu-T)w+\bar m w^2].
\]
Choose the normal width small enough that the bracket in the second
equation is less than $-2\beta_*$, and that the variation of $T$ before
normal contraction is contained in the larger $T$-interval. Extend the
field smoothly outside this larger interval, preserving $T=0$, $w=0$
and this contraction. One can do this by bounded cutoffs in the
coefficient functions multiplying their displayed factors. The extension
is used only to define a complete auxiliary flow.

Define its fast projection by
\[
 \pi_0(T,w)=\lim_{t\to\infty}T(t).
\]
The decay $|w(t)|\le C|w|e^{-2\beta_*t}$ makes $\dot T$ integrable.
The variational system has a negative normal diagonal, an integrable
horizontal diagonal, and horizontal and coefficient couplings containing
$w(t)$. Variation of constants bounds horizontal variations and gives
normal decay. Higher normal forcings and derivatives of $\dot T$ retain
a normal state or lower normal variation; induction bounds them by
$C(1+t)^Me^{-\beta_*t}$, including frozen coefficient derivatives. Hence
\begin{equation}\label{rw:transitions:zero-phase}
 \pi_0(T,0)=T,\qquad \pi_0(0,w)=0,\qquad
 \pi_0-T=O\bigl(Tw(|T|+|w|)\bigr),\qquad
 \partial_T\pi_0>1/2.
\end{equation}
The order follows by integrating $Tw(\bar m w-T)$, whose factors the
flow preserves. A smaller normal width gives the derivative margin.

Use $\chi=\pi_0(T,w)$ as base coordinate. Its velocity vanishes
identically when $\epsilon=0$, and on $\chi=0$. Successive Hadamard
integration therefore writes it as $\epsilon\chi A$. Its time-one
state derivative in the base direction is $1+O(\epsilon)$, whereas
the normal multiplier is at most $e^{-\beta_*}+O(\epsilon)$.
The time-one graph transform is consequently a contraction on a small
set of graphs over the extended base. For its $k$th state derivative
the coefficient of the highest graph derivative is bounded by
\[
 e^{-\beta_*}(1+C|\epsilon|)^k+C|\epsilon|<1
\]
once $k$ is fixed and the scale cutoff is small enough. The lower
terms depend on lower graph derivatives. This proves convergence in
$C^k$ by induction, applied also to differences of consecutive iterates.

Ordinary $\epsilon$ differentiation introduces an $O(1)$ base
shear.  Order the induction first by total order, then by the number
of coefficient derivatives.  The highest unknown retains the same
contracting multiplier, and forcing of the same total order involves
fewer coefficient derivatives.  The scalar
variational equations give
\begin{equation}\label{rw:transitions:slow-variations}
 |D^\alpha\phi_t|\le C_\alpha(1+|t|)^{m_\alpha}
                       e^{C_\alpha|\epsilon||t|}.
\end{equation}
An ordinary scale derivative adds the time factor in the
variation-of-constants integral; iteration gives the displayed bound.
Normal contraction absorbs all finite polynomial factors. The graph
transform therefore converges with the prescribed mixed smoothness,
and uniqueness makes its graph invariant under the whole flow.

Return to the original $T$ parametrization. The graph is zero at
$\nu=0$ and at $\epsilon=0$. Its first ordinary $\epsilon$
derivative is also zero at $\epsilon=0$: differentiation of the
invariance equation has zero forcing there, because the independent
forcing in \eqref{rw:transitions:compact-lower} begins with
$\epsilon^2\nu T^2$. Finally its value and first $T$ derivative at
$T=0$ vanish. The point $(0,0)$ is invariant, and its weak and normal
linear eigendirections are distinct. These statements give the five
factors $\epsilon^2\nu T^2$ by finite-order Hadamard integration.
To retain order $N+2$ afterwards, construct the graph with at least
$N+7$ joint derivatives. Comparing the $T^2$ coefficient in its
invariance equation gives
\[
 2\epsilon\bar B\,w_{g,2}
   =\epsilon^2\nu+(-1+\epsilon\bar B)w_{g,2},
\]
which proves the exact value of $H(0)$.

Subtract the graph and reparametrize the base so that its coordinate
$x_0$ agrees with $T$ on the graph. In these coordinates
\[
 \dot x_0=\epsilon a(x_0)+\epsilon Z\Gamma(x_0,Z),
 \qquad \dot Z=-Z\widetilde\beta(x_0,Z).
\]
The factor $\epsilon$ survives because the zero-scale coordinate is
$\pi_0$; invariance of $T=0$ also makes $\Gamma(0,Z)=0$.
Let $\phi_t$ be the complete scalar flow of $\dot x_0=\epsilon a(x_0)$.
Set
\[
 x=\lim_{t\to\infty}\phi_{-t}(x_0(t)).
\]
Indeed,
\[
 \frac{\dd}{\dd t}\phi_{-t}(x_0(t))
   =\epsilon D\phi_{-t}(x_0(t))Z(t)\Gamma(x_0(t),Z(t)),
\]
and each selected derivative is bounded by
$C(1+t)^M e^{-(\beta_*-C_N|\epsilon|)t}$, by
\eqref{rw:transitions:slow-variations}. Reduce the scale cutoff so
that $C_N|\epsilon|<\beta_*/8$ for all the finitely many variational
products used below; the majorants are then integrable. The limit is a
coordinate with the prescribed finite smoothness, equals $x_0$ on
the graph, and differs from
$x_0$ by $\epsilon x_0ZO_{C^N}(1)$. The flow identity for the limit
gives $\dot x=\epsilon a(x)$ exactly. Restriction of the original
$T$ equation to the graph now yields
\eqref{rw:transitions:lower-factor}.

The normal linear eigenvalue at $T=w=0$ is exactly
$-1+\epsilon\bar B$. At zero scale its value along $w=0$ is $-1$.
Continuity, on the common compact domains and with the stated finite
jets, proves the normal estimates. Convexity follows from
\eqref{rw:transitions:lower-factor}. Finally $x_R=1/R$ and the
choice of $R$ give
$\bar B-\bar ax_R+x_R^2>3x_R^2/4$; reducing the scale cutoff
preserves a fixed positive margin there.
\end{proof}

\begin{lemma}[Sections and the actual exit]\label{rw:transitions:lower-sections}
On the fixed horizontal section $v=v_h>0$, with $v_h$ sufficiently
small, the phase input is
$t=x(T,v_h+\epsilon T)=T k(T)$, where $k>c>0$ and $t_T>c>0$.
The normal input has a positive lower bound. On the section $Y=0$
with $T$ in a fixed compact subset of $(0,T_R)$, it is instead
$w=\epsilon T$, the normal input is $\epsilon$ times a positive
bounded unit, and $x=T+O_{C^N}(\epsilon)$.

Restrict either entry to a fixed positive phase distance from $x_R$.
Its forward orbit reaches the fixed phase exit $x=x_R$ if and only if
\begin{equation}\label{rw:transitions:lower-passing}
g_I(s)>0\qquad(t\le s\le x_R).
\end{equation}
This criterion concerns the phase exit in the signed tube. Positivity
of the original central height is imposed separately after composition.
The entire transit remains in the chosen strip, and its duration
$T_f$ is at least $c/\epsilon$. Its selected finite time jets obey
\[
 |D^\alpha T_f|\le C\epsilon^{-M}(1+T_f)^M
                         e^{C_N\epsilon T_f}.
\]
\end{lemma}
\begin{proof}
Equation~\eqref{rw:transitions:zero-phase} and
$w_g=O(\epsilon^2T^2)$ give $Z>v_h/2$ on the horizontal entry and
positive-unit $Z/\epsilon$ on the finite entry. They also give the
phase margins, reversed in $q=1/T$ on fixed finite compact intervals.

The normal coordinate contracts, so it stays in the selected tube.
The scalar phase crosses exactly under \eqref{rw:transitions:lower-passing}:
an equilibrium blocks finite-time crossing, whereas positivity on the
closed interval gives a finite transit integral. The larger domain keeps
the original orbit in the unchanged field. Bounded $g_I$ and separated
sections give $T_f\ge c/\epsilon$. For the time jets apply the implicit
hitting calculation of Appendix~\ref{rw:app:relative-variations} to
$\phi_{T_f}(t)=x_R$, using \eqref{rw:transitions:slow-variations} and
the exit speed $\epsilon g_I(x_R)\ge c\epsilon$. Thus no minimum-gap
denominator is introduced.
\end{proof}
